\documentclass[11pt]{article}
% Shared preamble for the printable mini lessons (1.4 onward).
%   \documentclass[11pt]{article}
%   % Shared preamble for the printable mini lessons (1.4 onward).
%   \documentclass[11pt]{article}
%   % Shared preamble for the printable mini lessons (1.4 onward).
%   \documentclass[11pt]{article}
%   \input{../preamble}
% Build the key with:  pdflatex -jobname=<name>_KEY "\def\iskey{}\input{<name>}"
% Figures shared with the web page live in figures/Unit_X/<name>.tex and are
% pulled in with \fig{Unit_X/<name>}; the standalone package drops their own
% preamble, so their macros come from figures/preamble.tex, loaded here.
\input{../../figures/preamble}
\usepackage[margin=0.75in]{geometry}
\usepackage{array}
\usepackage{enumitem}
\usepackage{amsmath}
\usepackage[helvet]{sfmath}  % math in the same sans font as the text
\usepackage{standalone}
\pagestyle{empty}
\setlength{\parindent}{0pt}
\setlist[enumerate,1]{leftmargin=1.5em,itemsep=1em}
\setlist[enumerate,2]{label=(\alph*),leftmargin=2em,itemsep=2.2em}

\newif\ifkey
\ifdefined\iskey \keytrue \fi
\colorlet{keyred}{red!75!black}
\newcommand{\ans}[1]{\ifkey\leavevmode{\color{keyred}#1}\fi}
% vertical work space on the student copy, less on the key
\newcommand{\work}[1]{\ifkey\vspace{0.6em}\else\vspace{#1}\fi}
\newcommand{\fig}[1]{\input{../../figures/#1}}
% Title line: \lessontitle{Quartiles, IQR, and Box Plots}
\newcommand{\lessontitle}[1]{{\Large\bfseries #1\ifkey\ \textcolor{keyred}{KEY}\fi}\hfill Name \underline{\hspace{2in}}\par\vspace{0.6em}}
\newcommand{\partner}{\par\vspace{0.4em}Work with a partner. Compare answers when you finish each problem.\par}

% Build the key with:  pdflatex -jobname=<name>_KEY "\def\iskey{}\input{<name>}"
% Figures shared with the web page live in figures/Unit_X/<name>.tex and are
% pulled in with \fig{Unit_X/<name>}; the standalone package drops their own
% preamble, so their macros come from figures/preamble.tex, loaded here.
% Shared preamble for web figures.  Each figure is a standalone .tex that
% \input's this file; scripts/make-figures.sh turns it into an SVG.
\usepackage[T1]{fontenc}
\usepackage{helvet}
\renewcommand{\familydefault}{\sfdefault}
\usepackage{tikz}
\usetikzlibrary{decorations.pathreplacing}
% axis numbers in the text font, not math italic
\pgfkeys{/pgf/number format/.cd,assume math mode=true,1000 sep={}}

% dot plot: \dotplot{min}{max}{value/count,...}{axis label}
\newcommand{\dotplot}[4]{%
\begin{tikzpicture}[x=0.9cm,y=0.33cm]
  \draw[thick] (#1-0.5,0) -- (#2+0.5,0);
  \foreach \t in {#1,...,#2} {\draw (\t,0.25) -- (\t,-0.25) node[below] {\t};}
  \foreach \v/\c in {#3} {\ifnum\c>0 \foreach \k in {1,...,\c} {\fill (\v,\k*1.15) circle (0.16cm);}\fi}
  \node at ({(#1+#2)/2},-2.3) {#4};
\end{tikzpicture}}

% shape sketch histogram (12 bins): \shapehist{heights}{label}
\newcommand{\shapehist}[2]{%
  \foreach \h [count=\i from 0] in {#1} \draw[fill=gray!30] (\i,0) rectangle (\i+1,\h);
  \draw (-0.3,0) -- (12.3,0);
  \node[font=\bfseries] at (6,-2.6) {#2};}
\newcommand{\meanmark}[1]{\fill (#1,0) -- ++(-0.3,-0.75) -- ++(0.6,0) -- cycle;}
\newcommand{\medmark}[2]{\draw[very thick,dashed] (#1,0) -- (#1,#2);}

% number line with minor ticks: \nline{min}{max}{label step}{axis label}
\newcommand{\nline}[4]{%
  \draw[thick] (#1-0.5,0) -- (#2+0.5,0);
  \foreach \t in {#1,...,#2} \draw (\t,0.12) -- (\t,-0.12);
  \pgfmathtruncatemacro{\nxt}{#1+#3}
  \foreach \t in {#1,\nxt,...,#2} \draw (\t,0.2) -- (\t,-0.2) node[below] {\t};
  \node at ({(#1+#2)/2},-0.95) {#4};}
% box plot above the line: \bplot{min}{Q1}{med}{Q3}{max}
\newcommand{\bplot}[5]{%
  \draw[thick] (#1,1.0) -- (#2,1.0);
  \draw[thick] (#4,1.0) -- (#5,1.0);
  \draw[thick] (#1,0.75) -- (#1,1.25);
  \draw[thick] (#5,0.75) -- (#5,1.25);
  \draw[thick,fill=gray!20] (#2,0.55) rectangle (#4,1.45);
  \draw[very thick] (#3,0.55) -- (#3,1.45);}

% ---- added for 1.4 onward ----
% axis on y=0: \axis{min}{max}{label step}{tick step}{axis label}
\newcommand{\axis}[5]{%
  \draw[thick] (#1,0) -- (#2,0);
  \pgfmathsetmacro{\axT}{#1+#4}\foreach \t in {#1,\axT,...,#2} \draw (\t,0.1) -- (\t,-0.1);
  \pgfmathsetmacro{\axL}{#1+#3}\foreach \t in {#1,\axL,...,#2}
    \draw (\t,0.18) -- (\t,-0.18) node[below] {\pgfmathprintnumber{\t}};
  \node at ({(#1+#2)/2},-0.95) {#5};}
% box plot centred at height y: \bplotat{y}{min}{Q1}{med}{Q3}{max}
\newcommand{\bplotat}[6]{%
  \draw[thick] (#2,#1) -- (#3,#1);
  \draw[thick] (#5,#1) -- (#6,#1);
  \draw[thick] (#2,#1-0.22) -- (#2,#1+0.22);
  \draw[thick] (#6,#1-0.22) -- (#6,#1+0.22);
  \draw[thick,fill=gray!20] (#3,#1-0.42) rectangle (#5,#1+0.42);
  \draw[very thick] (#4,#1-0.42) -- (#4,#1+0.42);}
% an outlier dot: \outlier{x}{y}
\newcommand{\outlier}[2]{\fill (#1,#2) circle (2.4pt);}

\usepackage[margin=0.75in]{geometry}
\usepackage{array}
\usepackage{enumitem}
\usepackage{amsmath}
\usepackage[helvet]{sfmath}  % math in the same sans font as the text
\usepackage{standalone}
\pagestyle{empty}
\setlength{\parindent}{0pt}
\setlist[enumerate,1]{leftmargin=1.5em,itemsep=1em}
\setlist[enumerate,2]{label=(\alph*),leftmargin=2em,itemsep=2.2em}

\newif\ifkey
\ifdefined\iskey \keytrue \fi
\colorlet{keyred}{red!75!black}
\newcommand{\ans}[1]{\ifkey\leavevmode{\color{keyred}#1}\fi}
% vertical work space on the student copy, less on the key
\newcommand{\work}[1]{\ifkey\vspace{0.6em}\else\vspace{#1}\fi}
\newcommand{\fig}[1]{\input{../../figures/#1}}
% Title line: \lessontitle{Quartiles, IQR, and Box Plots}
\newcommand{\lessontitle}[1]{{\Large\bfseries #1\ifkey\ \textcolor{keyred}{KEY}\fi}\hfill Name \underline{\hspace{2in}}\par\vspace{0.6em}}
\newcommand{\partner}{\par\vspace{0.4em}Work with a partner. Compare answers when you finish each problem.\par}

% Build the key with:  pdflatex -jobname=<name>_KEY "\def\iskey{}\input{<name>}"
% Figures shared with the web page live in figures/Unit_X/<name>.tex and are
% pulled in with \fig{Unit_X/<name>}; the standalone package drops their own
% preamble, so their macros come from figures/preamble.tex, loaded here.
% Shared preamble for web figures.  Each figure is a standalone .tex that
% \input's this file; scripts/make-figures.sh turns it into an SVG.
\usepackage[T1]{fontenc}
\usepackage{helvet}
\renewcommand{\familydefault}{\sfdefault}
\usepackage{tikz}
\usetikzlibrary{decorations.pathreplacing}
% axis numbers in the text font, not math italic
\pgfkeys{/pgf/number format/.cd,assume math mode=true,1000 sep={}}

% dot plot: \dotplot{min}{max}{value/count,...}{axis label}
\newcommand{\dotplot}[4]{%
\begin{tikzpicture}[x=0.9cm,y=0.33cm]
  \draw[thick] (#1-0.5,0) -- (#2+0.5,0);
  \foreach \t in {#1,...,#2} {\draw (\t,0.25) -- (\t,-0.25) node[below] {\t};}
  \foreach \v/\c in {#3} {\ifnum\c>0 \foreach \k in {1,...,\c} {\fill (\v,\k*1.15) circle (0.16cm);}\fi}
  \node at ({(#1+#2)/2},-2.3) {#4};
\end{tikzpicture}}

% shape sketch histogram (12 bins): \shapehist{heights}{label}
\newcommand{\shapehist}[2]{%
  \foreach \h [count=\i from 0] in {#1} \draw[fill=gray!30] (\i,0) rectangle (\i+1,\h);
  \draw (-0.3,0) -- (12.3,0);
  \node[font=\bfseries] at (6,-2.6) {#2};}
\newcommand{\meanmark}[1]{\fill (#1,0) -- ++(-0.3,-0.75) -- ++(0.6,0) -- cycle;}
\newcommand{\medmark}[2]{\draw[very thick,dashed] (#1,0) -- (#1,#2);}

% number line with minor ticks: \nline{min}{max}{label step}{axis label}
\newcommand{\nline}[4]{%
  \draw[thick] (#1-0.5,0) -- (#2+0.5,0);
  \foreach \t in {#1,...,#2} \draw (\t,0.12) -- (\t,-0.12);
  \pgfmathtruncatemacro{\nxt}{#1+#3}
  \foreach \t in {#1,\nxt,...,#2} \draw (\t,0.2) -- (\t,-0.2) node[below] {\t};
  \node at ({(#1+#2)/2},-0.95) {#4};}
% box plot above the line: \bplot{min}{Q1}{med}{Q3}{max}
\newcommand{\bplot}[5]{%
  \draw[thick] (#1,1.0) -- (#2,1.0);
  \draw[thick] (#4,1.0) -- (#5,1.0);
  \draw[thick] (#1,0.75) -- (#1,1.25);
  \draw[thick] (#5,0.75) -- (#5,1.25);
  \draw[thick,fill=gray!20] (#2,0.55) rectangle (#4,1.45);
  \draw[very thick] (#3,0.55) -- (#3,1.45);}

% ---- added for 1.4 onward ----
% axis on y=0: \axis{min}{max}{label step}{tick step}{axis label}
\newcommand{\axis}[5]{%
  \draw[thick] (#1,0) -- (#2,0);
  \pgfmathsetmacro{\axT}{#1+#4}\foreach \t in {#1,\axT,...,#2} \draw (\t,0.1) -- (\t,-0.1);
  \pgfmathsetmacro{\axL}{#1+#3}\foreach \t in {#1,\axL,...,#2}
    \draw (\t,0.18) -- (\t,-0.18) node[below] {\pgfmathprintnumber{\t}};
  \node at ({(#1+#2)/2},-0.95) {#5};}
% box plot centred at height y: \bplotat{y}{min}{Q1}{med}{Q3}{max}
\newcommand{\bplotat}[6]{%
  \draw[thick] (#2,#1) -- (#3,#1);
  \draw[thick] (#5,#1) -- (#6,#1);
  \draw[thick] (#2,#1-0.22) -- (#2,#1+0.22);
  \draw[thick] (#6,#1-0.22) -- (#6,#1+0.22);
  \draw[thick,fill=gray!20] (#3,#1-0.42) rectangle (#5,#1+0.42);
  \draw[very thick] (#4,#1-0.42) -- (#4,#1+0.42);}
% an outlier dot: \outlier{x}{y}
\newcommand{\outlier}[2]{\fill (#1,#2) circle (2.4pt);}

\usepackage[margin=0.75in]{geometry}
\usepackage{array}
\usepackage{enumitem}
\usepackage{amsmath}
\usepackage[helvet]{sfmath}  % math in the same sans font as the text
\usepackage{standalone}
\pagestyle{empty}
\setlength{\parindent}{0pt}
\setlist[enumerate,1]{leftmargin=1.5em,itemsep=1em}
\setlist[enumerate,2]{label=(\alph*),leftmargin=2em,itemsep=2.2em}

\newif\ifkey
\ifdefined\iskey \keytrue \fi
\colorlet{keyred}{red!75!black}
\newcommand{\ans}[1]{\ifkey\leavevmode{\color{keyred}#1}\fi}
% vertical work space on the student copy, less on the key
\newcommand{\work}[1]{\ifkey\vspace{0.6em}\else\vspace{#1}\fi}
\newcommand{\fig}[1]{\input{../../figures/#1}}
% Title line: \lessontitle{Quartiles, IQR, and Box Plots}
\newcommand{\lessontitle}[1]{{\Large\bfseries #1\ifkey\ \textcolor{keyred}{KEY}\fi}\hfill Name \underline{\hspace{2in}}\par\vspace{0.6em}}
\newcommand{\partner}{\par\vspace{0.4em}Work with a partner. Compare answers when you finish each problem.\par}

% one row of an SD table: \sdrow{x}{x - mean}{(x - mean)^2}; the last two show on the key only
\newcommand{\sdrow}[3]{#1 & \ans{#2} & \ans{#3} \\ \hline}
\newenvironment{sdtable}{\renewcommand{\arraystretch}{1.45}\begin{tabular}{|c|p{1.1in}|p{1.1in}|}\hline
  $x$ & $x - \bar{x}$ & $(x - \bar{x})^2$ \\ \hline}{\end{tabular}}

\begin{document}

\lessontitle{Standard Deviation}

Two basketball players both score \textbf{10 points} per game on average.

\begin{center}\fig{Unit_2/2.1_players}\end{center}

Player A's points are close to the mean. Player B's points are spread far from the mean.
The \textbf{standard deviation} (SD) measures this.

\vspace{0.4em}
\begin{tabular}{|>{\raggedright\arraybackslash}p{6.7in}|}\hline
\textbf{Standard deviation} is about how far the values usually are from the mean.\newline
\textbf{Small SD:} values are close to the mean. \qquad \textbf{Large SD:} values are spread out.\\ \hline
\end{tabular}

\vspace{0.8em}
\textbf{How to find the standard deviation} \hfill \textbf{Example:} 7, 7, 10, 13, 13

\vspace{0.3em}
\begin{minipage}[t]{0.52\linewidth}\vspace{0pt}
\renewcommand{\arraystretch}{1.35}
\begin{tabular}{@{}l@{\;\;}l}
1. & Find the mean, $\bar{x}$. \quad $\bar{x} = 50 \div 5 = 10$\\
2. & Subtract the mean from each value.\\
3. & Square each result.\\
4. & Add the squares. \quad $9 + 9 + 0 + 9 + 9 = 36$\\
5. & Divide by $n - 1$. \quad $36 \div 4 = 9$\\
   & (This number is called the \textbf{variance}.)\\
6. & Take the square root. \quad $s = \sqrt{9} = \mathbf{3}$\\
\end{tabular}
\end{minipage}\hfill
\begin{minipage}[t]{0.42\linewidth}\vspace{0pt}\centering
\renewcommand{\arraystretch}{1.3}
\begin{tabular}{|c|c|c|}\hline
$x$ & $x - \bar{x}$ & $(x - \bar{x})^2$ \\ \hline
7 & $-3$ & 9 \\
7 & $-3$ & 9 \\
10 & 0 & 0 \\
13 & 3 & 9 \\
13 & 3 & 9 \\ \hline
\multicolumn{2}{|r|}{Sum} & 36 \\ \hline
\end{tabular}
\end{minipage}

\vspace{0.8em}
\textbf{The formula:} \qquad
$\displaystyle s = \sqrt{\frac{\text{sum of } (x - \bar{x})^2}{n - 1}}$ \qquad ($n$ = how many values)

\vspace{0.6em}
\textbf{Why $n - 1$?} Our data is usually a \textbf{sample}, a small part of a bigger group.
A sample is usually a little less spread out than the whole group. Dividing by $n - 1$ makes $s$ a little bigger to fix this.

\vspace{0.6em}
\textbf{Remember:}
\begin{itemize}[itemsep=0pt,topsep=2pt]
  \item The SD is never negative.
  \item The SD is 0 only when every value is the same.
  \item The SD has the same units as the data (points, minutes, inches).
\end{itemize}

\partner

% ---------------- PROBLEMS ----------------
\newpage
\begin{enumerate}

\item Find the standard deviation of 4, 6, 8, 10, 12.

\vspace{0.3em}
\begin{minipage}[t]{0.45\linewidth}\vspace{0pt}
\begin{sdtable}
\sdrow{4}{$-4$}{16}
\sdrow{6}{$-2$}{4}
\sdrow{8}{0}{0}
\sdrow{10}{2}{4}
\sdrow{12}{4}{16}
\multicolumn{2}{|r|}{Sum} & \ans{40} \\ \hline
\end{sdtable}
\end{minipage}\hfill
\begin{minipage}[t]{0.5\linewidth}\vspace{0pt}
Mean: \ans{\; $40 \div 5 = 8$}

\vspace{4em}
$s =$ \ans{\; $\sqrt{40 \div 4} = \sqrt{10} \approx 3.16$}
\end{minipage}
\work{2em}

\item Four quiz scores: 7, 3, 7, 7. Find the standard deviation.

\vspace{0.3em}
\begin{minipage}[t]{0.45\linewidth}\vspace{0pt}
\begin{sdtable}
\sdrow{7}{1}{1}
\sdrow{3}{$-3$}{9}
\sdrow{7}{1}{1}
\sdrow{7}{1}{1}
\multicolumn{2}{|r|}{Sum} & \ans{12} \\ \hline
\end{sdtable}
\end{minipage}\hfill
\begin{minipage}[t]{0.5\linewidth}\vspace{0pt}
Mean: \ans{\; $24 \div 4 = 6$}

\vspace{4em}
$s =$ \ans{\; $\sqrt{12 \div 3} = \sqrt{4} = 2$}
\end{minipage}
\work{2em}

\item Look at the two players on page 1.

\vspace{0.3em}
\begin{tabular}{@{}ll}
Player A: & 9, 10, 10, 11, 10\\
Player B: & 2, 18, 5, 15, 10\\
\end{tabular}
\begin{enumerate}[itemsep=1.2em]
  \item Find the mean of each player. \ans{\quad A: $50 \div 5 = 10$. \; B: $50 \div 5 = 10$.}
  \item Without calculating, which player has the larger SD? Explain.

  \ans{Player B. B's points are much farther from 10.}
  \work{1.5em}
  \item Find $s$ for Player A.

  \ans{Differences $-1, 0, 0, 1, 0$. Squares add to 2. \; $s = \sqrt{2 \div 4} = \sqrt{0.5} \approx 0.71$}
  \work{5em}
  \item Player B's SD is about 6.67. A coach needs a player who scores close to 10 points every game. Which player should the coach choose? Explain.

  \ans{Player A. A's SD is much smaller, so A's scores stay close to 10.}
\end{enumerate}

\newpage
\item All three dot plots have a mean of 5. Put them in order from the \textbf{smallest} SD to the \textbf{largest} SD. Explain.
\begin{center}\fig{Unit_2/2.1_rank_sd}\end{center}
\ans{R, P, Q. The dots in R are packed close to 5. The dots in Q are far from 5.\\
(The SDs are about 0.67, 1.83, and 3.16.)}
\work{4em}

\item Think about what the SD means.
\begin{enumerate}[itemsep=2.4em]
  \item Can a standard deviation be negative? Explain.

  \ans{No. It comes from adding squares (never negative) and taking a square root.}
  \item Find the SD of 10, 10, 10, 10 without a formula. Explain.

  \ans{0. Every value equals the mean, so there is no spread.}
  \item Set B is 9, 10, 10, 11. Set C is 0, 10, 10, 20. Both have a mean of 10. Which set has the larger SD? Explain.

  \ans{Set C. Its values are much farther from 10.}
\end{enumerate}

\end{enumerate}

\end{document}
